% Options for packages loaded elsewhere
\PassOptionsToPackage{unicode}{hyperref}
\PassOptionsToPackage{hyphens}{url}
%
\documentclass[
]{article}
\usepackage{amsmath,amssymb}
\usepackage{iftex}
\ifPDFTeX
  \usepackage[T1]{fontenc}
  \usepackage[utf8]{inputenc}
  \usepackage{textcomp} % provide euro and other symbols
\else % if luatex or xetex
  \usepackage{unicode-math} % this also loads fontspec
  \defaultfontfeatures{Scale=MatchLowercase}
  \defaultfontfeatures[\rmfamily]{Ligatures=TeX,Scale=1}
\fi
\usepackage{lmodern}
\ifPDFTeX\else
  % xetex/luatex font selection
\fi
% Use upquote if available, for straight quotes in verbatim environments
\IfFileExists{upquote.sty}{\usepackage{upquote}}{}
\IfFileExists{microtype.sty}{% use microtype if available
  \usepackage[]{microtype}
  \UseMicrotypeSet[protrusion]{basicmath} % disable protrusion for tt fonts
}{}
\makeatletter
\@ifundefined{KOMAClassName}{% if non-KOMA class
  \IfFileExists{parskip.sty}{%
    \usepackage{parskip}
  }{% else
    \setlength{\parindent}{0pt}
    \setlength{\parskip}{6pt plus 2pt minus 1pt}}
}{% if KOMA class
  \KOMAoptions{parskip=half}}
\makeatother
\usepackage{xcolor}
\setlength{\emergencystretch}{3em} % prevent overfull lines
\providecommand{\tightlist}{%
  \setlength{\itemsep}{0pt}\setlength{\parskip}{0pt}}
\setcounter{secnumdepth}{-\maxdimen} % remove section numbering
\ifLuaTeX
  \usepackage{selnolig}  % disable illegal ligatures
\fi
\usepackage{bookmark}
\IfFileExists{xurl.sty}{\usepackage{xurl}}{} % add URL line breaks if available
\urlstyle{same}
\hypersetup{
  pdftitle={Applied Focus: Imperfect-Information AI Systems},
  hidelinks,
  pdfcreator={LaTeX via pandoc}}

\title{Applied Focus: Imperfect-Information AI Systems}
\author{}
\date{}

\begin{document}
\maketitle

\section{Applied Focus: Imperfect-Information AI
Systems}\label{applied-focus-imperfect-information-ai-systems}

\subsection{Public Edition}\label{public-edition}

\textbf{Author:} Mauricio Garcia Villanueva\\
\textbf{Project:} Liar's Dice CFR Lab\\
\textbf{Live demo:} emgarviii.github.io/liars\_dice/

\subsection{Purpose}\label{purpose}

This applied focus paper connects the Liar's Dice CFR Lab to the broader
field of imperfect-information game solving. The original course
deliverable explored concepts like Nash equilibrium, minimax reasoning,
Counterfactual Regret Minimization, poker AI, and multi-agent learning.
This public edition keeps that direction, but cleans up the structure
and aligns the claims with the current validated demo.

The core idea is simple: many important decisions happen when agents do
not see the full state of the world. In games, that means hidden cards
or dice. In applied systems, it can mean incomplete market information,
uncertain counterparties, adversarial incentives, or delayed feedback.
The project uses a simplified dice game as a small, inspectable
environment for those ideas.

\subsection{Imperfect Information}\label{imperfect-information}

An imperfect-information game is a game where at least one player must
act without seeing the complete state. Poker is the standard example
because each player has private cards. Liar's Dice has the same kind of
structure because each player has private dice.

The important part is not the theme of the game. The important part is
the decision problem:

\begin{itemize}
\tightlist
\item
  What do I know?
\item
  What does the other player know?
\item
  What can the other player infer?
\item
  Which action is best when the missing information changes the value of
  every choice?
\end{itemize}

This is why imperfect-information games are useful for AI research. They
force an agent to reason over beliefs and incentives, not just visible
board positions.

\subsection{Multi-Agent Learning}\label{multi-agent-learning}

Multi-agent learning is different from ordinary supervised learning
because the environment includes other decision makers. If one agent
changes strategy, the other agent's best response may change too. That
creates a moving target.

In the Liar's Dice lab, the AI is not learning from a fixed dataset of
labeled examples. It trains through simulated interaction. The claimant
policy learns which claims are useful under hidden information. The
responder policy learns when to believe and when to challenge.

That self-play framing is one reason the project maps well to AI/ML
interviews. It shows the difference between fitting a static dataset and
designing a decision system inside an adversarial environment.

\subsection{Game-Solving Concepts}\label{game-solving-concepts}

Nash equilibrium is a strategy profile where no player can improve by
unilaterally changing strategy. It is a central concept in game theory,
but it is also easy to overclaim. A project can use sampled CFR+ methods
without proving that the final policy is a Nash equilibrium.

Minimax reasoning is another core idea. In zero-sum settings, each
player can think in terms of maximizing their own payoff while the
opponent tries to minimize it. That framing is powerful, but many
practical games and multi-agent systems are more complicated than a
clean two-player zero-sum example.

Counterfactual Regret Minimization is one of the most important
algorithmic ideas for extensive-form imperfect-information games. CFR
repeatedly updates a policy by comparing chosen actions against
alternative actions that could have been taken at the same information
set. Over time, regret minimization can produce strong average
strategies in the right settings.

The public demo uses sampled CFR+ training on a simplified game
abstraction. That wording is intentional. The current artifact
demonstrates the pipeline and measured behavior, but it does not claim a
formal equilibrium proof.

\subsection{Why Poker AI Matters}\label{why-poker-ai-matters}

Poker AI is the best-known applied example of imperfect-information game
solving. Systems like Libratus and Pluribus showed that self-play,
abstraction, and search can produce strategies strong enough to compete
with expert humans. Those systems also made clear that practical game
solving is not just one algorithm. It is a combination of abstraction,
offline training, evaluation, and real-time or targeted refinement.

The Liar's Dice project is much smaller, but the structure is similar in
spirit:

\begin{itemize}
\tightlist
\item
  reduce the game to a tractable representation,
\item
  train policies offline,
\item
  evaluate against opponent classes,
\item
  inspect where the policy fails,
\item
  explain what the result does and does not prove.
\end{itemize}

That structure is the applied lesson. I am not claiming that this
project is comparable in scale to poker AI systems. I am showing that I
understand the engineering pattern at a smaller scale and can build a
public artifact around it.

\subsection{Why Liar's Dice Works as a Teaching
Environment}\label{why-liars-dice-works-as-a-teaching-environment}

Liar's Dice is easy to explain. Each player rolls dice privately. A
claim is made about the total number of dice showing a face. The
responder has to decide whether the claim is likely enough to believe.

That makes it useful for communicating hidden information. A visitor can
immediately see their own dice and the hidden AI dice. If the claim is
``three of face four,'' the visitor has to combine visible evidence with
uncertainty about the opponent's hand.

The simplified public mode removes repeated raising so the decision is
easier to isolate. That choice makes the project more understandable as
an AI/ML demo, even though it means the solver is not solving the full
classic game.

\subsection{Applied Relevance}\label{applied-relevance}

The project has relevance beyond games because the same decision pattern
appears in other domains:

\begin{itemize}
\tightlist
\item
  A trader acts without seeing every other participant's position.
\item
  A fraud or risk system must infer intent from partial evidence.
\item
  A negotiation agent must reason about hidden preferences.
\item
  A market-making system faces adversarial selection and uncertainty.
\item
  A security model must perform well against adaptive behavior, not just
  fixed tests.
\end{itemize}

The Liar's Dice lab is not a finance model and should not be presented
as one. The connection is conceptual: hidden state, strategic agents,
probabilistic policies, and adversarial evaluation.

That is why I think the project fits AI/ML and quant-style
conversations. It demonstrates the kind of reasoning those areas
require, but in a compact environment where the rules and metrics are
visible.

\subsection{Evaluation as the Main Applied
Lesson}\label{evaluation-as-the-main-applied-lesson}

The strongest part of the public project is the evaluation story. The
updated policy conditions on remaining dice counts and private dice,
improving benchmark average win rate from 45.3 percent to 73.0 percent
across seeded opponent profiles. That is useful evidence that the policy
became better against the benchmark suite.

But benchmark success is not the same as robustness. The
best-response-style diagnostic still found exploitable pressure,
especially on the claimant side. That matters because in adversarial
systems, a weak opponent profile can make a policy look better than it
really is.

The public results page makes this distinction explicit. It explains
what the win rates mean, what opponents they are against, and why they
do not prove equilibrium.

\subsection{Public Framing}\label{public-framing}

The right way to describe the project is:

This is an end-to-end imperfect-information game-solving lab built
around a simplified Liar's Dice challenge game. It includes Python
sampled CFR+ self-play, policy export, deterministic evaluation, a
playable TypeScript demo, a classic rules comparison, and a curated
paper archive.

The wrong way to describe it would be to present it as a classic-game
solver.

Another overclaim would be to present it as mathematically certified
optimal play.

The public edition avoids those overclaims because the measured evidence
is more interesting than hype. The policy improved, but it remains
exploitable. That is a real research result.

\subsection{Takeaways}\label{takeaways}

Imperfect-information games are useful because they force agents to act
under hidden state.

Multi-agent learning is useful because the opponent is part of the
environment.

Sampled CFR+ training is useful because regret minimization gives a
principled way to update policies over repeated self-play.

Evaluation is necessary because benchmark win rates can hide
exploitability.

The Liar's Dice CFR Lab is valuable as a portfolio artifact because it
connects these ideas into one public, reproducible system.

\end{document}
